\section{Direct calibration of the abstention rule}
\label{sec:calibration}

The design document fixed a detectable halving of per-cell output ($s=0.5$),
coverage of at least $90\%$, and discrimination of at least $80\%$, defined as
the probability that the interval excludes zero at that effect. The proposed
calibration chooses the smallest tested count meeting both criteria for an
unqualified estimate, and the smallest meeting coverage alone for a
wide-interval estimate. The panel, labelling axes and falsification rule are
frozen in configuration and checked against exact values in the test suite.
These safeguards record the specification; they do not establish its adequacy.

\paragraph{The draw pool changes the result.}
The generator filters eligibility by held-out patient, pooling cells from both
tissues before drawing either synthetic arm. The design document did not
specify tissue eligibility. We therefore evaluate both this implemented pool
and a reference-tissue-only pool. With pooled cells, no count bin meets both
criteria: discrimination reaches at most $0.78$ across 13 seeds at 50
replicates, and $0.75$ across eight seeds at 500. The reference-only pool does
return candidates, but does not recover the committed count of 50
(Figure~\ref{fig:calibration}).

\begin{figure}[t]
\centering
\includegraphics[width=\linewidth]{figures/fig1_calibration.pdf}
\caption{Direct coverage and discrimination evaluation of the cell-count rule.
Dashed lines mark the pre-specified targets; shading marks the committed
three-state rule. Left: the implemented pooled-cell generator returns no
qualifying cutpoint at any tested count. Right: the reference-only generator
returns candidates whose reported values depend on the grid ($100$, $90$ and
$70$ across three grids), rather than the committed 50. These candidates do
not carry patient-level uncertainty estimates.}
\label{fig:calibration}
\end{figure}

\paragraph{Reported cutpoints depend on binning and grid resolution.}
The committed grid varies mature-cell fractions at 2{,}000 cells per arm,
reaching counts $\{0,20,40,100,200,400,800\}$. It has no point between 40 and
100. On the reference pool it returns $ok=100$ on all eight seeds at 500
replicates, and twelve of thirteen at 50. A superseded report gave 120, which
this configuration cannot produce; the provenance defect is recorded in
Appendix~\ref{app:withdrawn}.

Coverage and discrimination are aggregated into twelve log-spaced count bins.
Each reported count is the median of the counts landing in its bin, rather
than a directly estimated crossing. Adding nine fractions to the grid moves
$ok$ from 100 to 90 while retaining the same selected bin, $(76,107]$:
the committed grid contributes $\{100\}$ there, while the filled grid
contributes $\{80,90,100\}$. The additional counts also change the evidence
aggregated within the bin. The newly populated bin $(55,76]$ has median
discrimination $0.789$, below target.

A third grid, reaching below 20, returns $ok=70$ on all eight seeds. Its count
summaries are $\{2.5,10,17.5,27.5,42.5,70,105,150,400,800\}$; fractional
values are summaries, not fractional cells. Discrimination is similar near
shared count ranges: medians over eight seeds are $0.712$, $0.728$ and $0.726$
near 40, and $0.880$, $0.858$ and $0.878$ near 100. The reported crossing
brackets and candidates nevertheless differ:

\begin{center}\small
\begin{tabular}{@{}lccc@{}}
\toprule
 & \textbf{committed} & \textbf{filled} & \textbf{dense} \\
\midrule
crossing bracketed to & $(40,100]$ & $(65,90]$ & $(42.5,70]$ \\
reported $ok$ & $100$ & $90$ & $70$ \\
reported $wide$ & $40$ & $45$ & $42$ \\
\bottomrule
\end{tabular}
\end{center}

These brackets describe adjacent reported count summaries, not confidence
intervals. The decreasing candidates across these three grids establish
sensitivity to the tested configurations, not monotonicity under further
refinement. No tested grid, cohort or seed returned the committed 50.
The $wide$ candidate has zero seed-to-seed variance in all six cohort-by-grid
combinations ($40$, $45$, $42$ here; $100$, $65$, $70$ in the second cohort).
That stability alone cannot establish precision when only a few output values
are available. Pooled draws return no qualifying cutpoint on any grid or seed
in either cohort; the dense-grid discrimination maxima are $0.750$ and $0.795$.

\paragraph{Simulation precision is not patient-level uncertainty.}
On the filled reference grid, 50 replicates yield candidates between 45 and 90
across seeds; at 500, seven of eight seeds return 90. More replicates reduce
Monte Carlo noise conditional on the available cohort. Ten patients with two
held out per replicate provide only 45 distinct holdout pairs, so 500
replicates revisit each pair about eleven times on average. This does not
quantify uncertainty from sampling patients from a larger population.
We therefore report grid-dependent candidates, not an endorsed calibrated
cutpoint. Appendix~\ref{app:gridcaveats} records denominator and coincident-cutpoint
issues, and Appendix~\ref{app:guards} examines small-sample interval calibration.

\paragraph{A second cohort gives the same calibration failure.}
The second Lee cohort contains six patients, compared with ten in the first.
On the committed and filled grids, pooled draws return no qualifying cutpoint
at either replicate count. Discrimination alone reaches $0.80$ at 50
replicates but at most $0.66$ at 500. Reference-only draws yield $ok=400$ or
800, eight to sixteen times the committed 50. At 50 replicates, eleven of
thirteen seeds return 800; at 500, six of eight return 400. This cohort has
only fifteen distinct holdout pairs, and its largest grid point requests 800
mature cells per arm from a pool of 844 distinct mature cells.

Thus neither the recovery curve nor seed stability justified the committed
rule. The algebraic equality replicates across both cohorts, while direct
calibration remains sensitive to pool, grid and cohort. On real samples,
abstention dominates at the finest annotation resolution
(Appendix~\ref{app:modal}); the calibration failure does not license replacing
those unavailable contrasts with zeros.
